\section*{Chapter \ref{ch:iv1:mental-technique} --- Mental Technique}
\markboth{\protect\omTocMark{}{Answers: Mental Technique}}{\protect\omTocMark{}{Answers: Mental Technique}}

\begin{solution}{iq:iv1:mental-technique:1}
The product is 9991.

Use the difference of squares\index{difference of squares}. The two numbers sit 3 either side of 100, so $97 \times 103 = (100 - 3)(100 + 3) = 100^2 - 3^2 = 9991$.

The idea works whenever the two numbers have an even sum and their average is easy to square. For $47 \times 53$ the average is 50, which gives $2500 - 9 = 2491$. For $96 \times 108$ the average is 102 and each number is 6 away from it, which gives $10\,404 - 36 = 10\,368$.

\iqfollow{And if the sum is odd, as in $47 \times 54$?}{Peel off one copy first: $47 \times 54 = 47 \times 53 + 47 = 2491 + 47 = 2538$.}
\iqwrong{Saying ``about 10\,000''. That is a good size check, but an interviewer asking for the product wants the 9 taken off. The correction is a single small square, so it costs nothing.}
\end{solution}

\begin{solution}{iq:iv1:mental-technique:2}
$85^2 = 7225$. The rule: multiply the leading part by one more than itself, and then write 25.

For 85 the leading part is 8, and $8 \times 9 = 72$, so the answer is 72 followed by 25. To prove it, write the number as $10a + 5$:
\[
(10a + 5)^2 = 100a^2 + 100a + 25 = 100\,a(a + 1) + 25.
\]
The rule holds for any $a$, not only one digit: $115^2$ has $a = 11$, and $11 \times 12 = 132$ gives $13\,225$.

\iqfollow{Is there a similar rule near 50?}{Yes. $(50 + d)^2 = 2500 + 100d + d^2$, so add $d$ to 25 for the hundreds and write $d^2$ in the last two places. $53^2 = 2809$, and with $d = -3$, $47^2 = 2209$.}
\iqfollow{How do you square any other two-digit number?}{Move to the nearest round number and correct: $n^2 = (n + d)(n - d) + d^2$. So $67^2 = 70 \times 64 + 9 = 4489$.}
\end{solution}

\begin{solution}{iq:iv1:mental-technique:3}
The product is 9312. Take 3 from 96, or equally 4 from 97, to get the 93. Then multiply the distances from 100, $3 \times 4 = 12$, and write it in the last two places.

To see why, let the distances be $a$ and $b$:
\[
(100 - a)(100 - b) = 100\,(100 - a - b) + ab.
\]
The first term is the left part of the answer, and $ab$ is the right part. The same algebra handles numbers above 100, where the distances add: $103 \times 106 = 100 \times 109 + 18 = 10\,918$.

\iqwrong{Treating the right part as two digits no matter what. If $ab$ reaches 100 it carries. For $88 \times 85$ the distances are 12 and 15, so the left part is 73 and $ab = 180$, and the answer is $7300 + 180 = 7480$.}
\iqfollow{Does this need 100 as the base?}{No. Any easy base works. Near 50, $48 \times 47 = 50 \times 45 + 6 = 2256$.}
\end{solution}

\begin{solution}{iq:iv1:mental-technique:4}
The product is $46\,000$.

Since $125 = 1000/8$, multiply by 1000 and divide by 8. In practice do the division first, because it keeps the numbers small: $368 / 8 = 46$, and then attach three zeros.

The trick is to replace a hard multiplier by an easy fraction of a power of ten. The set worth knowing is $5 = 10/2$, $25 = 100/4$, $125 = 1000/8$, $12.5 = 100/8$ and $75 = 300/4$. If the division does not come out whole, keep going in decimals: $372 / 8 = 46.5$, so $372 \times 125 = 46\,500$.

\iqfollow{What about $368 \times 375$?}{$375 = 3 \times 125$, so three times the answer above: $138\,000$.}
\iqfollow{And dividing by 125?}{Do the opposite: multiply by 8 and divide by 1000. $3500 / 125 = 28\,000 / 1000 = 28$.}
\end{solution}

\begin{solution}{iq:iv1:mental-technique:5}
The product is 5346.

Write the first digit, then the sum of each pair of neighbours, then the last digit, working from the right so that you can carry. For 486 the columns are $6$, then $8 + 6 = 14$, then $4 + 8 = 12$, then $4$. Going from the right: write 6; write 4 and carry 1; $12 + 1 = 13$, so write 3 and carry 1; $4 + 1 = 5$. The result is 5346.

It works because $11n = 10n + n$. Writing $n$ shifted one place above $n$ itself puts each digit in a column with its neighbour, and column addition then gives exactly those sums:
\[
4860 + 486 = 5346.
\]

\iqwrong{Forgetting the carry and running the column sums together, which gives something like 4\,12\,14\,6. Doing the two-line sum is safer until the carries are automatic.}
\iqfollow{What is the matching trick for 101?}{Writing the number twice with a two-place shift: $486 \times 101 = 48\,600 + 486 = 49\,086$.}
\end{solution}

\begin{solution}{iq:iv1:mental-technique:6}
$16\%$ of 25 is 4. The same holds for 25\% of 16, which is what you actually compute.

The swap is legal because both are the same product over 100: $p\%$ of $n$ is $pn/100$, and multiplication does not care about the order. So choose whichever way round has an easy percentage. Thus 18\% of 50 is 50\% of 18, which is 9, and 4\% of 175 is 175\% of 4, which is 7.

\iqwrong{Treating this as a party trick that only works with round numbers. It always holds. It only helps when one of the two numbers is an easy percentage: 25, 50, 75, 125, or a multiple of 10. For 37\% of 41, neither way round is easy.}
\iqfollow{What is 36\% of 25, and 12.5\% of 64?}{36\% of 25 is 25\% of 36, which is 9. And 12.5\% is one eighth, so 12.5\% of 64 is 8.}
\end{solution}

\begin{solution}{iq:iv1:mental-technique:7}
The tax is 64, and the pre-tax price is 320.

The price including tax is 120\% of the pre-tax price, so the tax is $20/120$ of it, which is one sixth, and $384 / 6 = 64$. In general, a tax of rate $r$ on the net makes up $r/(1 + r)$ of the gross.

\iqwrong{Taking 20\% of 384, which gives 76.8. The 20\% applies to the pre-tax price, not to the amount that already includes the tax.}
\iqfollow{What share of the gross is a 25\% markup, and an 8\% one?}{A 25\% markup is one fifth of the gross, since $25/125 = 1/5$. An 8\% markup is $8/108 = 2/27$, or about $7.4\%$.}
\iqfollow{You paid 384 after a 20\% discount. What was the list price?}{You paid 80\% of the list price, so it was $384 / 0.8 = 480$. The discount of 96 is a quarter of what you paid, and that is the same reverse step.}
\end{solution}

\begin{solution}{iq:iv1:mental-technique:8}
A rise of 25\% is needed. In general, after a fall of $d$, the rise needed is $d/(1 - d)$. A fall of 37.5\% needs a rise of 60\%.

After the fall you hold $1 - d$ of the start, and the rise must bring it back to 1, so it must be $\frac{1}{1 - d} - 1 = \frac{d}{1 - d}$. Fractions make this quick: a fall of one fifth leaves four fifths, and you need to gain one quarter. A fall of $37.5\% = 3/8$ leaves $5/8$, and you need to gain $3/5$. A fall of one half needs a rise of 100\%.

\iqwrong{Saying 20\%. The rise is taken on a smaller base, so it must be larger in percentage terms.}
\iqwhy{A strategy that is 20\% below its high must earn 25\% to recover, and one that is 50\% below must double. So drawdowns cost more than they look, and the cost grows quickly with their size.}
\end{solution}

\begin{solution}{iq:iv1:mental-technique:9}
$11/14$ is larger, by $1/126$.

Cross-multiply\index{cross-multiplication}: $a/b$ and $c/d$ compare in the same way as $ad$ and $bc$, as long as the denominators are positive. Here $7 \times 14 = 98$ and $9 \times 11 = 99$, so the second is larger. The difference is the gap in cross products over the product of denominators:
\[
\frac{11}{14} - \frac{7}{9} = \frac{99 - 98}{14 \times 9} = \frac{1}{126}.
\]
A second route compares the distances from 1: $2/9 \approx 0.222$ and $3/14 \approx 0.214$, and the smaller gap belongs to the larger fraction.

\iqwrong{Saying the fraction with the larger numbers must be larger. It says nothing: $1/2$ exceeds $11/30$.}
\iqfollow{Which is larger, $13/17$ or $23/30$?}{$13 \times 30 = 390$ and $17 \times 23 = 391$, so $23/30$ is larger, by $1/510$.}
\end{solution}

\begin{solution}{iq:iv1:mental-technique:10}
$3/7 = 0.428571\dots$ and $5/7 = 0.714285\dots$. Keep the same six digits in the same cyclic order, and start from whichever digit gives the right size.

For $3/7$, you know it is about $0.43$, so the cycle must start at 4: 428571. For $5/7$, about $0.71$, it starts at 7: 714285.

The reason is long division. Dividing 1 by 7 passes through the remainders $1, 3, 2, 6, 4, 5$ and back to 1, which is every possible nonzero remainder. Starting the division from 3 instead of 1 only means joining that same loop at a different point, so you read the same digits from a different place.

\iqfollow{Does every fraction behave this way?}{No. It needs the remainders to form one loop. For 13 the remainders split into two loops, so $1/13 = 0.076923\dots$ and $3/13 = 0.230769\dots$ share a cycle, but $2/13 = 0.153846\dots$ does not.}
\end{solution}

\begin{solution}{iq:iv1:mental-technique:11}
The average price is 85.

Use deviations from a base\index{weighted average by deviations}. Take 80 as the base. The 500 shares at 88 stand 8 above it, and they make up $5/8$ of the total, so the average sits $5/8$ of the way from 80 to 88. That is $80 + 5 = 85$.

A weighted average of two numbers always lies between them, at a distance proportional to the other side's weight. So you need only the ratio of the weights, never the totals.

\iqfollow{You hold 400 at $101.5$, 100 at 102 and 500 at 101. What is your average?}{Use a base of 101. The deviations are $0.5$, 1 and 0, with weights $0.4$, $0.1$ and $0.5$. Their weighted sum is $0.2 + 0.1 = 0.3$, so the average is $101.3$.}
\iqwrong{Averaging the prices without weights: $(80 + 88)/2 = 84$. The larger order pulls the average towards its own price, so the answer must be above 84.}
\end{solution}

\begin{solution}{iq:iv1:mental-technique:12}
The rule gives 9 years, and the exact answer is $9.006$ years. It is that close at 8\%.

Solve $(1 + r)^t = 2$, which gives $t = \ln 2 / \ln(1 + r)$. For small $r$, $\ln(1 + r) \approx r$, so $t \approx 0.693/r$, or $69.3$ divided by the rate in percent. Then $\ln(1 + r)$ is a little below $r$, which pushes the numerator up: at 8\% the exact numerator is about $72.05$. The number 72 is also chosen because it divides by 2, 3, 4, 6, 8, 9 and 12.

\iqwrong{Using simple interest: $100/8 = 12.5$ years. That ignores interest on the interest, which is the whole point of doubling rules.}
\iqfollow{How long does it take to triple at 8\%?}{Use $\ln 3 \approx 1.1$ in the same way. At this rate the numerator rises to about 114, so the answer is about $114/8 = 14.25$ years. The exact value is $14.27$.}
\end{solution}

\begin{solution}{iq:iv1:mental-technique:13}
About 2\% a day. The quick number is 16.

Variances of independent returns add, so the standard deviation grows like the square root of time\index{square-root-of-time rule}. Annual volatility is therefore daily volatility times $\sqrt{252} \approx 15.87$. Using 16 is about 1\% off, which is far better than any estimate of volatility itself. So $32\% / 16 = 2\%$; the exact figure is $2.02\%$.

The other two numbers worth knowing are $\sqrt{52} \approx 7.2$ for weeks and $\sqrt{12} \approx 3.46$ for months. At 32\% annual volatility a week's return has a standard deviation of about $4.4\%$, and a month's about $9.2\%$.

\iqwrong{Dividing by 252, which gives $0.13\%$. That treats volatility as if it added up over time as a return does. Variance adds, not volatility.}
\iqwhy{This is the reflex behind every daily risk figure. Seeing a stock move 4\% in a day at 32\% annual volatility, you know at once that it is a two-standard-deviation day.}
\end{solution}

\begin{solution}{iq:iv1:mental-technique:14}
$\log_{10} 72 \approx 1.857$ and $\log_{10} 0.0072 \approx -2.143$.

Factor the number into primes you know, and turn the product into a sum\index{logarithms by factoring}. Since $72 = 2^3 \times 3^2$,
\[
\log_{10} 72 = 3(0.301) + 2(0.477) = 0.903 + 0.954 = 1.857.
\]
The exact value is $1.8573$. For the second, $0.0072 = 72 \times 10^{-4}$, so subtract 4 and get $-2.143$.

You also have $\log_{10} 5$ for free: $5 = 10/2$, so it is $1 - 0.301 = 0.699$. These three, with powers of 10, cover every number made of the primes 2, 3 and 5. That is why the pair is worth memorising.

\iqfollow{What is $\log_{10} 1.5$?}{$1.5 = 3/2$, so it is $0.477 - 0.301 = 0.176$.}
\iqwrong{Trying to take the log of $0.0072$ as a positive number between 0 and 1. A number below 1 has a negative log, so a sign check catches this at once.}
\end{solution}

\begin{solution}{iq:iv1:mental-technique:15}
$2^{30}$ is about $10^9$, a billion. One trillion takes about 40 doublings. The first estimate is about 7\% low.

The fact is that $2^{10} = 1024 \approx 10^3$. So $2^{30} = (2^{10})^3 \approx 10^9$, and $2^{40} \approx 10^{12}$. The true $2^{30}$ is $1\,073\,741\,824$, or $1.024^3 \approx 1.074$ billion. The error compounds at about $2.4\%$ per factor of 1024, so it matters more for high powers: $2^{100}$ is about $1.27 \times 10^{30}$ rather than $10^{30}$.

\iqfollow{How large is $2^{64}$?}{$2^{60} \approx 1.15 \times 10^{18}$, using $1.024^6 \approx 1.15$, and $2^4 = 16$, so about $1.8 \times 10^{19}$.}
\iqwhy{Developers read sizes this way all the time. A table with a billion keys needs about 30 bits per key to tell them apart, and it runs out of a 32-bit range quickly.}
\end{solution}

\begin{solution}{iq:iv1:mental-technique:16}
$\sqrt{50} \approx 7.071$. The estimate $7 + 1/14 = 7.0714$ is slightly too high, and the true value is $7.07107$.

Take the nearest easy square, $a^2 = 49$, with $d = 50 - 49 = 1$, and use
\[
\sqrt{a^2 + d} \approx a + \frac{d}{2a}.
\]
This is one step of Newton's method\index{Newton's method}. Squaring the estimate shows why it always overshoots: $(a + d/2a)^2 = a^2 + d + (d/2a)^2$, which is above the target by a positive amount, for either sign of $d$. So one more correction, always downward, brings it far closer, and it still stays just above the true value.

\iqfollow{Estimate $\sqrt{3\,200\,000}$.}{Write it as $1000 \sqrt{3.2}$. Since $1.8^2 = 3.24$, take $a = 1.8$ and $d = -0.04$. That gives $1.8 - 0.04/3.6 = 1.789$, so about 1789. The true value is $1788.9$.}
\iqwrong{Interpolating linearly between $\sqrt{49}$ and $\sqrt{64}$, which gives $7 + 1/15$. It uses a slope that is too flat for this part of the curve, and so it comes out low by about $0.004$.}
\end{solution}

\begin{solution}{iq:iv1:mental-technique:17}
10 basis points of \$3.5 million is \$3500, and the extra interest over three months is \$6250.

Remember one fixed point: one basis point of \$1 million is \$100. So $n$ basis points of $m$ million dollars is $100\,nm$ dollars: $100 \times 10 \times 3.5 = 3500$. For the loan, one basis point of \$250 million is \$25\,000 a year, and three months is a quarter of that, or \$6250.

\iqwrong{Confusing a basis point with $0.01$, which is a whole percent. As a decimal, one basis point is $0.0001$. If your answer for the loan looks like \$625\,000, you have made this error.}
\iqfollow{One eighth of a percent is how many basis points?}{12.5, since 1\% is 100 basis points.}
\iqwhy{Interest, spreads and fees are all quoted in basis points, and the conversion to dollars has to come without thinking.}
\end{solution}

\begin{solution}{iq:iv1:mental-technique:18}
99-17 is $99.53125$.

Know that $1/32 = 0.03125$, by halving down from one half: $0.5$, $0.25$, $0.125$, $0.0625$, $0.03125$. Then build the fraction from pieces: $17/32 = 16/32 + 1/32 = 0.5 + 0.03125 = 0.53125$.

\iqfollow{Going the other way, how would $101.40625$ be quoted?}{The fractional part is $0.40625$. Take away $0.375$, which is $12/32$, and $0.03125$ is left, which is $1/32$. So it is $13/32$, or 101-13.}
\iqfollow{What is the dollar value of a price move of one 32nd on \$10 million of face value?}{Price is per 100 of face, so a 32nd is $0.03125$ per 100, or $0.03125\%$ of face, which is \$3125.}
\iqwrong{Reading 99-17 as $99.17$. The digits after the hyphen count 32nds, not hundredths, so it is more than half a point above 99 and not about a sixth.}
\end{solution}

\begin{solution}{iq:iv1:mental-technique:19}
$38 \times 42$ is larger, by exactly 5.

Both pairs add to 80, so both have average 40. Write each as $40$ plus or minus its half-gap:
\[
37 \times 43 = 40^2 - 3^2 = 1591, \qquad 38 \times 42 = 40^2 - 2^2 = 1596.
\]
The general rule: with the sum fixed, the product is the square of the average minus the square of the half-gap\index{difference of squares}. So the closer the numbers are together, the larger their product. It is largest when they are equal, and each step apart costs the square of the half-gap, not a fixed amount.

\iqfollow{Two numbers add to 100. What is the largest possible product?}{Fifty each gives 2500. Gaps of 2, 10 and 20 give 2499, 2475 and 2400, so the loss is small near the top and quickly larger away from it.}
\iqfollow{Does the rule extend to more factors?}{Yes. With the sum fixed, the product of positive numbers is largest when they are equal. This is the inequality between the arithmetic and geometric means\index{arithmetic-geometric mean inequality}.}
\end{solution}

\begin{solution}{iq:iv1:mental-technique:20}
$1.02 \times 0.97 = 0.9894$. For a rough answer drop the cross term, which leaves $0.99$, with an error of $0.0006$.

Expand with $a = 0.02$ and $b = -0.03$:
\[
(1 + a)(1 + b) = 1 + a + b + ab = 1 + 0.02 - 0.03 - 0.0006 = 0.9894.
\]
The first-order part, $1 + a + b$, is the sum of the two moves. The cross term $ab$ is the product of two small numbers, so it is much smaller, and its sign is the product of their signs.

\iqwrong{Getting the sign of the cross term wrong and answering $0.9906$. With one move up and one down, the cross term is negative.}
\iqfollow{And $1.02 \times 1.03 \times 0.96$?}{First order: $1 + 0.02 + 0.03 - 0.04 = 1.01$. The pairwise products are $0.0006 - 0.0008 - 0.0012 = -0.0014$, and the triple product is about $-0.00002$. That gives $1.00858$, and the exact value is $1.008576$.}
\iqwhy{Returns over a few periods add to first order. This tells you how large a mistake that is.}
\end{solution}

\begin{solution}{iq:iv1:mental-technique:21}
It is wrong, because it fails the digit-sum check. It passes the checks on the last digit and on the size. The true product is $100\,022$.

Last digit: $7 \times 6 = 42$ ends in 2, and so does the claim. Size: $3850 \times 26 = 100\,100$ is very close to the claim. So neither helps.

Casting out nines\index{casting out nines}: a number has the same remainder modulo 9 as the sum of its digits, because $10 \equiv 1$ modulo 9. Here 3847 has digit sum 22, which reduces to 4. And 26 reduces to 8. So the product must reduce to $4 \times 8 = 32$, which reduces to 5. The claim's digits add to 6, not 5, so it is wrong.

\iqfollow{If the check passes, is the product right?}{No. A pass only says that the error is a multiple of 9, which a random error is about one time in nine. The check can prove a product wrong, but not right.}
\iqwhy{This takes a few seconds and catches most slips, including one wrong digit, which is the usual error. Use it on your own sums before you say them.}
\end{solution}

\begin{solution}{iq:iv1:mental-technique:22}
It is 4964, chosen by casting out nines\index{casting out nines}.

Every option ends in 4, as it must, since $3 \times 8 = 24$. So the last digit does not help. Size does not settle it easily either, since $70 \times 70 = 4900$ only tells you that the answer is in the region. So compare remainders modulo 9. 73 has digit sum 10, which reduces to 1, and 68 has digit sum 14, which reduces to 5. So the product must reduce to 5. The options reduce to $4964 \to 23 \to 5$, $4954 \to 22 \to 4$, $5064 \to 15 \to 6$ and $4864 \to 22 \to 4$. Only the first fits.

To check it directly, the difference of squares works with half-integers: the average is $70.5$ and the half-gap is $2.5$, so $70.5^2 - 2.5^2 = 4970.25 - 6.25 = 4964$.

\iqwrong{Choosing the option nearest to $70 \times 70 = 4900$. That picks 4864, which is wrong. A rounded product says where the answer lies, not which of several close options it is.}
\end{solution}

\begin{solution}{iq:iv1:mental-technique:23}
A digit-sum check cannot see this error because 6428 is 6248 with two digits swapped, and swapping digits leaves the sum alone. Use a check modulo 11, or the fact that 8 divides 88. The true product is 6248.

The digit sums agree: 71 reduces to 8, 88 reduces to 7, and $8 \times 7 = 56$ reduces to 2; and 6428 has digit sum 20, which reduces to 2. So nines pass.

For eleven, $10 \equiv -1$ modulo 11\index{casting out elevens}, so a number has the same remainder as the alternating sum of its digits from the right. Since $88 = 8 \times 11$, the product must leave remainder 0. But $8 - 2 + 4 - 6 = 4$, so 6428 fails. The true 6248 has $8 - 4 + 2 - 6 = 0$. The factor 8 also catches it, as $428 / 8$ is not whole, while $248 / 8 = 31$.

\iqfollow{Why does eleven catch swaps?}{Swapping neighbours moves each digit to a place of opposite sign, so the alternating sum changes by twice their difference, and that is not a multiple of 11 unless the digits are equal.}
\end{solution}

\begin{solution}{iq:iv1:mental-technique:24}
The net change is $+2\%$, and the order does not matter.

Multiply the factors: $1.2 \times 0.85 = 1.02$. Multiplication does not care about the order. The rule for combining two moves is $a + b + ab$:
\[
0.20 - 0.15 + (0.20)(-0.15) = 0.05 - 0.03 = 0.02.
\]
So add the two moves, and then add their product, watching its sign. The product is always of the two moves together, so it is negative when they have opposite signs.

\iqwrong{Saying $+5\%$, from adding the two percentages. From 100 the rise adds 20, and the fall, taken on the higher base of 120, removes 18 rather than 15.}
\iqfollow{And rising 10\% and then falling 10\%?}{$-1\%$: the moves cancel and only the cross term is left. In the same way, $+50\%$ and $-50\%$ leave you $25\%$ down.}
\end{solution}

\begin{solution}{iq:iv1:mental-technique:25}
$1/0.97 \approx 1.0309$ and $1/1.03 \approx 0.9709$.

Use the geometric series: for small $x$,
\[
\frac{1}{1 - x} = 1 + x + x^2 + \dots, \qquad \frac{1}{1 + x} = 1 - x + x^2 - \dots.
\]
With $x = 0.03$ the terms are $0.03$, $0.0009$ and $0.000027$, so $1/0.97 \approx 1.030927$ and $1/1.03 \approx 0.970873$. The next term is below $10^{-6}$, so three terms are good to about six decimal places.

The method stops working as $x$ grows: each term is only $x$ times the last, so for $x = 0.2$ you need many terms. Switch to fractions there: $1/0.8 = 5/4 = 1.25$.

\iqwrong{Saying $1/0.97 \approx 1.03$ and stopping. That is right to two decimal places, but the question asks for four, and the square term supplies them.}
\iqwhy{This is how to invert an exchange rate or a discount factor on the spot. A rate of $0.97$ in one direction is $1.0309$ in the other, not $1.03$. On a large order those missing $0.0009$ are 9 basis points.}
\end{solution}

\begin{solution}{iq:iv1:mental-technique:26}
$1.05^{10} \approx 1.63$. The exact value is $1.62889$.

First way, repeated squaring: $1.05^2 = 1.1025$. Squaring gives $1.05^4 \approx 1.2155$. Multiply by $1.05$ to get $1.05^5 \approx 1.2763$. Squaring that gives $1.629$. The squares look hard but are easy: $1.1025^2 \approx 1 + 2(0.1025) + 0.1025^2 \approx 1.2155$.

Second way, logs\index{exponential approximation}: $\ln 1.05 \approx 0.05 - 0.00125 + 0.00004 = 0.04879$, so ten of them give $0.4879$. Then $e^{0.4879} = e^{0.5} e^{-0.0121} \approx 1.6487 \times 0.988 \approx 1.629$.

\iqwrong{Saying $e^{0.5} \approx 1.65$, using $\ln(1 + x) \approx x$. Since $\ln(1 + x)$ is below $x$, this always comes out too high. The other wrong answer is the simple $1.50$, which ignores compounding.}
\iqfollow{What does the rule of 72 say?}{Doubling takes $72/5 = 14.4$ years, so ten years gives $2^{10/14.4} = 2^{0.694} \approx 1.62$. That is rougher, but it takes only seconds.}
\iqfollow{How can you check it?}{Use the binomial expansion: $1 + 10(0.05) + 45(0.0025) + 120(0.000125) + 210(0.00000625) \approx 1.6288$. The terms fall quickly, so five are enough.}
\end{solution}

\begin{solution}{iq:iv1:mental-technique:27}
$0.99^{100} \approx 0.366$. The exact value is $0.36603$.

Take logs\index{exponential approximation}: $\ln 0.99 = -0.01 - 0.00005 - \dots \approx -0.01005$. So $100 \ln 0.99 \approx -1.005$, and
\[
0.99^{100} \approx e^{-1} e^{-0.005} \approx 0.3679 \times 0.995 \approx 0.366.
\]

The answer $1/e \approx 0.368$ uses $\ln(1 - x) \approx -x$ and misses the next term of the series, $-x^2/2$. Over 100 periods the missing piece is $-100 \times 0.00005 = -0.005$, which is a half-percent change and shows up in the third decimal.

\iqfollow{What is $1.01^{100}$?}{The same way, $100 \ln 1.01 \approx 0.995$, so $e \times 0.995 \approx 2.705$. The exact value is $2.70481$. Note that $0.99^{100}$ and $1/1.01^{100} = 0.36971$ are different numbers, and both differ from $1/e$.}
\iqwhy{This is the chance of no event in 100 independent tries at 1\% each. So 100 tries at 1\% give at least one event with probability about $63.4\%$, not 100\%.}
\end{solution}

\begin{solution}{iq:iv1:mental-technique:28}
At $r = 1$ the rule gives 72 years against an exact $69.66$, which is $3.4\%$ too long. At $r = 36$ it gives 2 years against an exact $2.25$, which is $11\%$ too short. It is exact near $r = 7.85$. A good fix is a numerator of $69.3 + r/3$.

The exact doubling time is $\ln 2 / \ln(1 + r/100)$. Writing it as a numerator over $r$, the numerator is $r \ln 2 / \ln(1 + r/100)$, which is not constant: it rises with $r$ from $69.3$. Its values are about $69.66$ at $r = 1$, $72.05$ at $r = 8$ and $81.2$ at $r = 36$.

Expanding $\ln(1 + x)$ to second order gives $69.3\,(1 + x/2)$ with $x = r/100$, a slope of $0.35$ per percentage point. A slope of $1/3$ is easier to use. At $r = 36$ it gives $81.3/36 = 2.26$, and at $r = 1$ it gives $69.6$.

\iqwrong{Thinking that 72 is the theoretical constant. The theoretical constant is $69.3$. The 72 is a convenient number that happens to fit near typical rates.}
\end{solution}

\begin{solution}{iq:iv1:mental-technique:29}
A total return of about $28.6\%$. The simple answer is $25.2\%$, and it is too low because it ignores compounding.

Convert to a continuous rate\index{exponential approximation}. $\ln 1.001 \approx 0.001 - 0.0000005 = 0.0009995$, and $252$ of these give $0.2519$. Then
\[
e^{0.2519} \approx 1 + 0.2519 + \tfrac12 (0.2519)^2 + \tfrac16 (0.2519)^3 \approx 1.2863,
\]
and the next term of the series adds $0.00017$. The exact value of $1.001^{252}$ is $1.28643$, so the return is $28.6\%$.

\iqfollow{What if you lose 10 basis points a day instead?}{$0.999^{252} \approx e^{-0.2521} \approx 0.777$, a loss of $22.3\%$. So the gain and the loss are not of equal size, and the two together make $1.2864 \times 0.7771 < 1$.}
\iqfollow{Over five years, at the same daily rate?}{The rate in the exponent is five times larger, about $1.26$, so the factor is $e^{1.26} \approx 3.5$. Here the simple sum of $126\%$ is far too low.}
\end{solution}

\begin{solution}{iq:iv1:mental-technique:30}
About $60.5\%$ of your money is left, which is a loss of $39.5\%$.

A rise and a fall together multiply your money by $1.1 \times 0.9 = 0.99$, and there are 50 such pairs, so the result is $0.99^{50}$. Taking logs, $50 \ln 0.99 \approx -0.5025$, and $e^{-0.5025} \approx 0.605$.

The arithmetic average return is zero, but money compounds, and it compounds at the geometric rate\index{volatility drag}. The rule is that the growth rate per period is about the mean minus half the variance, $\mu - \sigma^2/2$. Here $\sigma = 0.1$, so the drag is $0.005$ per period, or $0.5$ over 100 periods, and $e^{-0.5} \approx 0.607$.

\iqwrong{Saying that you end where you started because the moves cancel. They cancel in addition, but returns multiply.}
\iqfollow{And at 20\% up and down?}{$0.96^{50} \approx 0.130$, compared with $e^{-2} \approx 0.135$ from the rule. The rule is getting rough, because the higher-order terms now matter.}
\iqwhy{Daily volatility of 1\% costs about $0.01^2/2 \times 252 = 1.26\%$ a year of compound growth. Double the volatility and the cost is four times larger, because it goes with the square.}
\end{solution}

\begin{solution}{iq:iv1:mental-technique:31}
About 23 years, and more exactly $22.76$. The rule of 72 gives 24, which is too long.

The exact time solves $0.97^t = 0.5$, so $t = \ln 2 / (-\ln 0.97)$. Expanding, $-\ln(1 - x) = x + x^2/2 + \dots$ is a little above $x$, so the denominator is about $0.0305$ and $t$ is below $69.3/3 = 23.1$.

Doubling goes the other way. There, $\ln(1 + x) = x - x^2/2 + \dots$ is a little below $x$, so the doubling time at 3\% is longer, at $23.45$ years. So at the same rate, growth takes longer to double than decay takes to halve. The rule of 72 is too long for both at 3\%.

\iqfollow{Is there a corrected numerator for decay?}{Yes. Where doubling uses $69.3 + r/3$, decay uses $69.3 - r/3$. At 3\% that is $68.3/3 = 22.8$ years.}
\iqwrong{Dividing 50\% by 3\%, which gives $16.7$ years. That applies 3\% of the original amount each year, whereas here the 3\% applies to what is left, and so the losses shrink.}
\end{solution}

\begin{solution}{iq:iv1:mental-technique:32}
$\log_{10} 7 \approx 0.845$. The true value is $0.8451$.

Look for a power of 7 close to a number made only of 2, 3 and 5\index{logarithms by factoring}. Since $7^4 = 2401$, and $2400 = 2^5 \times 3 \times 5^2$,
\[
\log_{10} 7 \approx \tfrac14 \log_{10} 2400 = \tfrac14 \bigl[ 5(0.301) + 0.477 + 2(0.699) \bigr] = \tfrac14 (3.380) = 0.845,
\]
using $\log_{10} 5 = 1 - 0.301 = 0.699$. The ratio $2401/2400$ differs from 1 by only $1/2400$, and dividing by four makes the error smaller still, so the rounding in the inputs is the only real error.

A rougher route uses $7^2 = 49 \approx 50$, which gives $(1.699)/2 = 0.8495$, wrong in the third decimal.

\iqwrong{Interpolating linearly between $\log_{10} 6 = 0.778$ and $\log_{10} 8 = 0.903$, which gives $0.8405$. The log curve bends, so a straight line through two points on it lies below the curve.}
\iqfollow{And $\log_{10} 13$?}{Use the same idea with $13 \times 5 = 65$, next to $64 = 2^6$. So $\log_{10} 13 = \log_{10} 64 + \log_{10}(65/64) - \log_{10} 5$. With $\log_{10}(65/64) \approx 0.0155 \times 0.434 \approx 0.0067$, this is $1.806 + 0.0067 - 0.699 \approx 1.114$.}
\end{solution}

\begin{solution}{iq:iv1:mental-technique:33}
$2^{100}$ has 31 digits, and its first digit is 1. It is about $1.27 \times 10^{30}$.

Take logs: $100 \times 0.30103 = 30.103$. The integer part says the number is between $10^{30}$ and $10^{31}$, so it has 31 digits. The fractional part gives the leading digits: $10^{0.103}$. Since $\log_{10} 1.25 = 1 - 3(0.30103) = 0.0969$, the remainder $0.103 - 0.0969 = 0.006$ multiplies by about $1.014$, which gives $1.25 \times 1.014 \approx 1.268$. The exact number starts $1.2676$.

\iqfollow{And for $3^{50}$?}{$50 \times 0.47712 = 23.856$, so 24 digits. Since $\log_{10} 7 \approx 0.845$, the leading part is a little above 7. It is $7.18 \times 10^{23}$.}
\iqwrong{Counting digits as the integer part of the log, here 30. A number with log $30.103$ has 31 digits, just as $10^{30}$ itself has 31.}
\iqwhy{Reading the size of a power is how you know at once whether a count fits a 64-bit integer. Since $2^{64}$ has 20 digits, a 21-digit count does not fit.}
\end{solution}

\begin{solution}{iq:iv1:mental-technique:34}
$\ln 1.3 \approx 0.262$. The true value is $0.26236$.

Find a nearby number whose log you know, and take a small step from it. Here $1.25 = 10/8$, so
\[
\ln 1.25 = \ln 10 - 3 \ln 2 = 2.303 - 2.079 = 0.224.
\]
Then $1.3 / 1.25 = 1.04$, and $\ln 1.04 \approx 0.04 - 0.0008 = 0.0392$. The total is $0.263$, and using more digits of the two given logs brings it to $0.262$.

The series $\ln(1 + x) = x - x^2/2 + x^3/3 - \dots$ with $x = 0.3$ is slow: the first three terms are $0.3$, $-0.045$ and $+0.009$. Stopping after two gives $0.255$ and after three $0.264$. A small step from a nearby known point needs only one or two terms.

\iqfollow{Find $\ln 3$ from these same two logs.}{Since $3^4 = 81$ is very near $80 = 2^3 \times 10$, you get $4 \ln 3 \approx 3(0.693) + 2.303 + \ln(81/80) = 4.382 + 0.0124$, so $\ln 3 \approx 1.099$.}
\iqwhy{A 30\% simple return is a log return of about $26.2\%$. The gap is why returns of that size cannot just be added up.}
\end{solution}

\begin{solution}{iq:iv1:mental-technique:35}
The tenth root of 2 is about $1.072$, and the fifth root of 3 is about $1.246$.

A root becomes a division in logs\index{exponential approximation}. For the first, $\ln 2 / 10 = 0.0693$, and then
\[
e^{0.0693} \approx 1 + 0.0693 + \tfrac12 (0.0693)^2 + \tfrac16 (0.0693)^3 \approx 1.0718.
\]
For the second, $\ln 3 / 5 = 0.2198$, and $1 + 0.2198 + 0.0242 + 0.0018 + 0.0001 \approx 1.246$. The true value is $1.2457$.

This ties to the rule of 72: the tenth root of 2 says doubling in ten years needs $7.18\%$ a year, and the rule gives $7.2\%$.

\iqwrong{Dividing the excess by 10 and saying $1.1$. That is simple interest, not compound. A growth rate of 10\% a year for ten years gives $2.59$, not 2.}
\iqfollow{What monthly rate compounds to 12\% a year?}{$\ln 1.12 \approx 0.1133$, and dividing by 12 gives $0.00944$, so about $0.949\%$ a month, not 1\%.}
\end{solution}

\begin{solution}{iq:iv1:mental-technique:36}
The product is about $4.47$ million, within $0.3\%$ of the exact value of $4\,484\,646$.

Round to numbers you can multiply: $50 \times 60 \times 80 \times 20 = 4\,800\,000$. Then correct using each rounding as a percentage. The rounded values are too high by $6.4\%$ ($50/47$) and $5.3\%$ ($20/19$), and too low by $3.2\%$ ($60/62$) and $1.2\%$ ($80/81$). For small errors these add, so the total is about $+7.3\%$. Divide $4.8$ million by $1.073$ to get about $4.47$ million.

Good rounding pairs one number up with another down, so that the corrections partly cancel and stay small enough to add. You have done this without a single long multiplication.

\iqwrong{Rounding all four and stopping at $4.8$ million. That is 7\% off, though it looks fine because every factor seems close. Four small errors add up, so they must be corrected.}
\iqfollow{What if the errors are not small?}{Multiply the ratios rather than adding them, or add logs. Adding percentages is only good while each stays within about 10\%.}
\end{solution}

\begin{solution}{iq:iv1:mental-technique:37}
The chance of exactly 10 heads is about $0.18$ (exactly $0.1762$), and of exactly 12 heads about $0.12$ (exactly $0.1201$).

Use the normal curve\index{normal approximation to the binomial}. The number of heads has mean 10 and variance $20/4 = 5$, so a standard deviation of $\sqrt 5 \approx 2.24$. The peak height of a normal density is $0.4$ divided by the standard deviation, which gives $0.4/2.24 \approx 0.18$. So the chance at the mean is close to that.

For 12 heads, go two away from the mean, and multiply by the normal shape $e^{-d^2/(2\sigma^2)} = e^{-4/10} \approx 0.67$. That gives $0.18 \times 0.67 \approx 0.12$.

\iqwrong{Saying one half because the coin is fair. One half is the share of heads you expect, not the chance of hitting exactly that count. There are 21 possible counts, and 10 is only the likeliest of them.}
\iqfollow{And for 100 tosses?}{The standard deviation is 5, so the chance of exactly 50 heads is about $0.4/5 = 0.08$, and the exact value is $0.0796$. Because it shrinks as the toss count grows, in proportion to one over its square root, you need four times the tosses to halve it.}
\end{solution}

\begin{solution}{iq:iv1:mental-technique:38}
$\binom{50}{25} \approx 1.3 \times 10^{14}$ (exactly $1.264 \times 10^{14}$), and $\binom{100}{30} \approx 2.9 \times 10^{25}$ (exactly $2.937 \times 10^{25}$).

For the central case, use the normal peak\index{normal approximation to the binomial}: the chance of exactly $n$ heads in $2n$ fair tosses is about $1/\sqrt{\pi n}$, which makes $\binom{2n}{n} \approx 4^n / \sqrt{\pi n}$. Here $4^{25} = 2^{50} \approx 1.126 \times 10^{15}$ using $1.024^5 \approx 1.126$, and $\sqrt{25\pi} \approx 8.86$, so the estimate is $1.27 \times 10^{14}$.

For an off-centre choice, use entropy\index{entropy estimate of a binomial coefficient}: with $p = 0.3$,
\[
\binom{N}{pN} \approx \frac{2^{N H(p)}}{\sqrt{2\pi N p (1 - p)}}, \qquad H(p) = -p \log_2 p - (1 - p) \log_2 (1 - p).
\]
Here $H(0.3) \approx 0.881$, so the numerator is $2^{88.1} \approx 3.3 \times 10^{26}$. The denominator is $\sqrt{2\pi \times 21} \approx 11.5$. This gives $2.9 \times 10^{25}$.

\iqwrong{Using $50^{25}/25!$, which is a good upper bound only when the selection is a small part of the whole. Here it is too large by a factor of about 1500.}
\iqfollow{And $\binom{52}{26}$?}{$\binom{52}{26} = \binom{50}{25} \times \frac{52 \times 51}{26 \times 26}$, a factor of $3.92$, so about $4.96 \times 10^{14}$.}
\end{solution}

\begin{solution}{iq:iv1:mental-technique:39}
About 226 bits, and exactly $225.58$ bits. $52!$ has 68 digits.

The number of orders is $52!$, so the fewest bits on average is $\log_2 52!$. Use Stirling's formula\index{Stirling's formula}, $\ln n! \approx n \ln n - n + \tfrac12 \ln(2\pi n)$, in base 2:
\[
\log_2 n! \approx n \log_2 n - \frac{n}{\ln 2} + \tfrac12 \log_2 (2\pi n).
\]
Since $52 = 32 \times 1.625$ and $\log_2 1.625 \approx 0.70$, you get $\log_2 52 \approx 5.70$. The terms are then $52 \times 5.70 = 296.4$, $52/0.693 = 75.0$ and $\tfrac12 \log_2 327 \approx 4.2$. The total is $296.4 - 75.0 + 4.2 = 225.6$.

For the digits, multiply by $\log_{10} 2$: $225.6 \times 0.301 \approx 67.9$, which means $52!$ is about $8 \times 10^{67}$, with 68 digits.

\iqwrong{Multiplying 52 by 6 bits a card, to get 312. That is a valid encoding but not the fewest: it allows a card to repeat. Removing that waste saves over 80 bits.}
\iqwhy{A 64-bit random seed can produce only $2^{64}$ shuffles out of $2^{226}$. If your simulations depend on the whole shuffle, a generator with a small seed reaches almost none of them.}
\end{solution}

\begin{solution}{iq:iv1:mental-technique:40}
The present value is about $\$12.4$, and the bond is worth about $75.2$ per 100 of face value.

An annuity\index{annuity factor} is a perpetuity minus a perpetuity that starts late. A perpetuity of \$1 a year is worth $1/r$ today, and one that starts after 30 years is worth that much at year 30, so $1.07^{-30}/r$ today. Subtracting,
\[
\text{PV} = \frac{1 - 1.07^{-30}}{0.07}.
\]
For the power, 7\% roughly doubles money every ten years: $1.07^{10} = 1.967$, so close to 2. Three of these give $1.07^{30} \approx 7.6$; the exact value is $7.61$. So $1/7.61 \approx 0.131$, and $(1 - 0.131)/0.07 = 0.869/0.07 \approx 12.4$. The exact value is $12.41$.

The bond pays 5 a year for 30 years and returns 100 at the end. So its price is $5 \times 12.41 + 100/7.61 = 62.05 + 13.14 = 75.2$.

\iqwrong{Using the middle payment date: $30 / 1.07^{15} \approx 10.9$. Discounting is convex, so early payments gain more than late ones lose, and this comes out too low.}
\iqfollow{What level yearly payment repays a loan of 100 over 30 years at 7\%?}{Divide by the factor: $100/12.41 = 8.06$ a year.}
\iqfollow{How much of a perpetuity's value is this 30-year annuity?}{A perpetuity is worth $1/0.07 = 14.29$, so the annuity is $1 - 1/7.61$, or about 87\%, of that. The payments after year 30 are worth only 13\%.}
\end{solution}
